% !TEX program = xelatex
% !TEX root = ../../TFM.tex
\documentclass[../../TFM.tex]{subfiles}

\begin{document}

    
    \chapter{Literatura empírica}
	\label{chap:literatura-empirica}
	\thesisstartchap{literatura-empirica}
	
	Este capítulo revisa la literatura empírica desde la lógica del modelo econométrico de la tesis. Más que presentar un inventario de estudios aislados, la discusión se organiza alrededor de los mecanismos mediante los cuales la dependencia de recursos naturales no renovables y la especialización en petróleo, gas y minería pueden relacionarse con la transformación estructural externa. Esta estrategia permite vincular la evidencia previa con dos resultados empleados en la investigación: el índice de complejidad económica (ECI), que aproxima la sofisticación de la canasta exportadora, y la diversificación exportadora ($DIVX=1-HHI$), definida como el complemento de la concentración de Herfindahl--Hirschman. La pertinencia de observar concentración y diversificación exportadora en economías intensivas en recursos se apoya en estudios que muestran que los recursos naturales pueden concentrar la estructura exportadora, desplazar exportaciones no extractivas y condicionar las posibilidades de diversificación \parencite{Harding2016,Bahar2018,Joya2015,Lashitew2021}. En consecuencia, la literatura se ordena en función de los resultados y canales incorporados en el modelo. Las referencias a GMM, cointegración, causalidad, ARDL u otros métodos describen las estrategias de los estudios citados y no implican que la tesis adopte esos procedimientos.
	
	\section{Medición empírica de la transformación estructural}
	\label{sec:medicion-transformacion-estructural}
	
	La literatura empírica sobre transformación estructural ha recurrido a indicadores que captan dimensiones distintas pero complementarias del cambio productivo. Una primera tradición se concentra en la dimensión externa de la estructura económica y utiliza medidas derivadas del patrón exportador para aproximar la acumulación de capacidades productivas. En esta línea, \textcite{Hidalgo2009} proponen un marco metodológico para inferir capacidades subyacentes a partir de redes país--producto construidas con datos de exportación. En su enfoque, la complejidad económica surge de la interacción entre la diversificación de los países y la ubicuidad de los productos, de modo que las economías capaces de exportar una variedad amplia de bienes poco ubicuos revelan un acervo más sofisticado de conocimientos y capacidades. El aporte central de este enfoque es haber convertido la estructura exportadora en una medida empírica de transformación productiva orientada hacia afuera.
	
	La utilidad del índice de complejidad económica\footnote{El Índice de Complejidad Económica (ECI) se interpreta en esta tesis como una aproximación indirecta a las capacidades productivas reveladas en la canasta exportadora. Por construcción, no captura plenamente capacidades productivas orientadas al mercado interno ni actividades no transadas internacionalmente.} se observa en la evidencia comparada internacional. Utilizando datos para 128 países, \textcite{Hausmann2014} muestran que el índice explica una fracción sustancial de la variación del ingreso per cápita, especialmente entre economías con baja dependencia de recursos naturales. Cuando una parte importante del ingreso proviene de rentas petroleras o mineras, la relación entre complejidad e ingreso se debilita, lo que indica que la abundancia de recursos puede elevar el ingreso sin traducirse necesariamente en una estructura productiva más compleja.
	
	La literatura también ha desarrollado extensiones y alternativas al índice de complejidad económica. Entre ellas, la métrica de \emph{fitness} propuesta por \textcite{Tacchella2012,Tacchella2013} introduce una relación no lineal entre países y productos y busca capturar con mayor precisión la estructura de capacidades productivas implícita en la matriz exportadora. Desde esta perspectiva, la \emph{fitness} de un país depende de su capacidad para sostener una canasta diversificada de bienes complejos, mientras que la complejidad de un producto se determina por el nivel de capacidades del país menos sofisticado que logra producirlo. Más allá de las diferencias metodológicas, tanto el índice de complejidad económica como la \emph{fitness} comparten la idea de que la transformación estructural puede observarse empíricamente en la calidad, diversidad y conectividad de la canasta exportadora.

	\textcite{Inoua2023} propone medir la complejidad a partir de una idea más directa: contar cuántos productos distintos exporta un país y tomar el logaritmo de esa cantidad. El autor sostiene que una canasta exportadora más diversa refleja una mayor variedad de conocimientos productivos, aunque advierte que las exportaciones de recursos naturales pueden distorsionar esta comparación cuando elevan los ingresos sin ampliar realmente la base productiva.
	
	Para los fines de esta tesis, la transformación estructural se observa desde la inserción externa de las economías. En ese plano, la complejidad económica y la diversificación exportadora captan dimensiones relacionadas pero no equivalentes: el ECI se concentra en la sofisticación y ubicuidad de los productos exportados, mientras que $DIVX=1-HHI$ permite observar si la canasta exportadora se distribuye entre un conjunto más amplio de productos. Esta distinción es especialmente importante en economías dependientes de recursos naturales, donde la literatura muestra que los ingresos extractivos pueden reforzar la concentración exportadora o desplazar exportaciones no asociadas a recursos \parencite{Harding2016,Bahar2018}.
	
	Las aplicaciones recientes del índice de complejidad económica confirman, además, que esta medida se asocia con resultados relevantes para la resiliencia y el desempeño estructural de los territorios. \textcite{Kazemzadeh2023} encuentran una relación no lineal entre complejidad e intensidad energética, mientras que \textcite{Rossouw2023} muestran que los municipios sudafricanos con menor complejidad presentan mayor vulnerabilidad frente a shocks económicos. Desde una perspectiva ambiental, \textcite{Aluko2023} muestran para 35 países de la OCDE durante 1998--2017 que el efecto de la complejidad económica sobre distintas medidas de degradación ambiental depende del nivel de ingreso. Para América Latina, \textcite{Alvarado2021} encuentran que los efectos de la complejidad económica y de las rentas de recursos naturales sobre la huella ecológica son heterogéneos a lo largo de la distribución, según estimaciones de regresión cuantílica.

	\section{Dependencia extractiva, concentración productiva y transformación estructural}
	\label{sec:dependencia-concentracion-transformacion}
	
	La evidencia empírica reciente muestra que la dependencia de recursos naturales constituye uno de los canales más importantes para explicar trayectorias divergentes de transformación estructural. Una parte de esta literatura analiza la relación en sentido inverso al enfoque clásico de la maldición de los recursos,\footnote{Mientras la literatura clásica de la maldición de los recursos suele analizar si la dependencia o abundancia de recursos afecta el crecimiento o la transformación productiva, esta línea empírica invierte la dirección analítica y examina si mayores capacidades productivas o mayor complejidad económica contribuyen a reducir la dependencia relativa de recursos naturales.} es decir, preguntándose en qué medida una mayor complejidad productiva permite reducir la dependencia extractiva. En este sentido, \textcite{Canh2020} encuentran, para una muestra global de economías, que mayores niveles de complejidad económica tienden a disminuir las rentas de recursos naturales, aunque este efecto no es homogéneo entre grupos de países y depende de la medida de complejidad utilizada. Sus resultados sugieren que la relación entre capacidades productivas y dependencia extractiva es potencialmente bidireccional y que la especialización en recursos puede obstaculizar la acumulación de capacidades asociadas con trayectorias de diversificación.

	Para África subsahariana, \textcite{Kelly2024} encuentran un resultado similar: la complejidad económica reduce el agotamiento de recursos naturales en un panel de 36 países durante 1997--2017, incluso al emplear medidas alternativas de complejidad y distintas estrategias de estimación.

	La relación entre complejidad y dependencia de recursos también presenta rasgos no lineales. \textcite{Ul-Durar2023}, para un panel de economías ricas en recursos durante el período 2000--2021, introducen una especificación cuadrática y estiman modelos de regresión cuantílica en panel para captar heterogeneidad a lo largo de la distribución de las rentas de recursos. Sus resultados muestran que la relación puede adoptar tanto una forma de U como de U invertida, dependiendo del tipo de recurso y del nivel de extracción. En particular, encuentran que en niveles bajos de extracción la complejidad económica puede coexistir con mayores rentas de recursos, mientras que en niveles altos contribuye a reducirlas. Este hallazgo sugiere que la dependencia extractiva no opera de forma lineal y que los efectos de los recursos sobre la transformación estructural pueden variar según la intensidad de la especialización y el punto en el que se encuentre cada economía dentro de su trayectoria productiva.

	\textcite{Zhang2023} analizan 123 países durante 1995--2018 mediante un modelo espacial de panel para distinguir efectos directos y derrames territoriales de las rentas de recursos sobre la complejidad económica. Sus resultados muestran que la abundancia total de recursos reduce la complejidad local, aunque los efectos varían por tipo de recurso: las rentas petroleras y mineras se asocian con menor ECI, mientras que las rentas de gas presentan un patrón más favorable. Además, encuentran que la calidad institucional y la apertura financiera atenúan los efectos adversos de las rentas extractivas, mientras que mayores flujos de inversión extranjera directa pueden intensificarlos.

	En la misma dirección, \textcite{Olaniyi-Odhiambo2025-2} incorporan una estructura de causalidad asimétrica para estudiar la relación entre complejidad económica y rentas de recursos naturales en Nigeria y Sudáfrica durante 1970--2021. Sus resultados no encuentran causalidad simétrica, pero sí una relación bidireccional asimétrica en Sudáfrica cuando se consideran los shocks positivos de ambas variables.
	
	La evidencia sobre países exportadores de recursos muestra que los ingresos extractivos tienden a desplazar exportaciones no extractivas y a reducir la amplitud de la canasta exportadora no asociada a recursos \parencite{Harding2016,Bahar2018}. A su vez, una estructura concentrada en pocos productos o actividades puede aumentar la exposición a shocks externos y limitar los eslabonamientos que sostienen la diversificación productiva \parencite{Bahar2018,Joya2015,Lashitew2021}.
	
	Estos estudios vinculan la dependencia extractiva y la concentración productiva con menores niveles de complejidad económica y con canastas exportadoras menos diversificadas \parencite{Canh2020,Bahar2018,Joya2015,Lashitew2021}.

	\section{Instituciones, gobernanza y diversificación productiva}
	\label{sec:instituciones-gobernanza-diversificacion}
	
	La literatura empírica sobre recursos naturales y transformación estructural ha identificado en las instituciones un canal decisivo para explicar por qué economías con dotaciones similares de recursos exhiben trayectorias productivas divergentes. \textcite{Ajide2022} muestran, mediante modelos dinámicos GMM para países africanos, que la debilidad institucional amplifica el efecto negativo de la dependencia de recursos naturales sobre la complejidad económica.
	
	En esta misma línea, \textcite{AvomNdoya2024} muestran que la estabilidad del país favorece la complejidad económica. A partir de un panel de 118 países para 1995--2018 estimado mediante \textit{System GMM}, los autores encuentran que entornos más estables facilitan la acumulación de capacidades productivas, especialmente a través de tres canales: inversión extranjera directa, desarrollo financiero y capital humano.
	
	En una dirección convergente, \textcite{Olaniyi-Odhiambo2025-3} analizan economías africanas ricas en recursos y encuentran que la complejidad económica reduce las rentas de recursos naturales, aunque este efecto depende de la calidad institucional y de la capacidad estatal para orientar esas rentas hacia actividades productivas más diversificadas.

	De forma complementaria, \textcite{Olaniyi-Odhiambo2025} estudian países africanos ricos en recursos durante 1995--2021 y muestran que la calidad institucional modifica la relación entre riqueza natural y complejidad económica. Aunque la riqueza de recursos y la calidad institucional presentan efectos directos positivos, la interacción entre ambas variables puede debilitar la complejidad cuando las instituciones no alcanzan un umbral suficiente. Este resultado sugiere que la capacidad de transformar rentas naturales en sofisticación productiva depende no solo de la presencia de instituciones, sino de su calidad efectiva para disciplinar la gestión de los recursos.
	
	\textcite{Owjimehr2024} estudian las 20 economías con mayor dependencia de recursos naturales durante 2000--2020 mediante un modelo PMG-ARDL. Sus resultados muestran que las rentas de recursos reducen la complejidad económica en el largo plazo, mientras que un mayor control de la corrupción favorece mejoras en ese indicador.

	En conjunto, estos estudios sugieren que la calidad institucional no opera solo como variable de control, sino como un mecanismo de mediación entre rentas extractivas, acumulación de capacidades y diversificación productiva \parencite{Ajide2022,Olaniyi-Odhiambo2025-3,Olaniyi-Odhiambo2025}.

	\textcite{Avom2022} analizan 108 países durante 1995--2017 mediante estimaciones GMM y encuentran que la abundancia de recursos naturales se asocia de manera robusta con menores niveles de complejidad económica. El efecto negativo es más claro en países en desarrollo y se mantiene al utilizar medidas alternativas de recursos y de complejidad. Además, los autores muestran que las instituciones políticas moderan esta relación, especialmente en democracias parlamentarias.

	\section{Capacidades productivas, innovación y nivel de desarrollo}
	\label{sec:capacidades-innovacion-desarrollo}
	
	La literatura empírica reciente coincide en que la transformación estructural depende no solo de la composición sectorial de una economía, sino también de sus dotaciones de capital humano, de su capacidad de innovación y de sostener niveles de ingreso compatibles con procesos de diversificación productiva. \textcite{YaseminYalta2021} analizan países de la región MENA mediante un modelo dinámico \textit{System GMM} y encuentran un efecto positivo del capital humano sobre la complejidad económica, junto con un impacto negativo de las rentas de recursos naturales.

	En una muestra amplia de países para 1995--2019, \textcite{Chu2023} confirma que el ingreso, la densidad poblacional, el capital humano y el gasto público se asocian positivamente con la complejidad económica, mientras que las rentas de recursos naturales tienen un efecto negativo. El uso de regresiones cuantílicas permite mostrar, además, que estos efectos varían según el nivel inicial de complejidad y el grupo de ingreso de cada país.

	Para el caso de China, \textcite{Khan2020} estudian la relación entre inversión extranjera directa y complejidad económica mediante modelos ARDL y VECM. Sus resultados muestran una relación bidireccional de largo plazo entre ambas variables, así como causalidad de corto plazo, lo que sugiere que la inversión extranjera puede operar como canal de transferencia de conocimiento, tecnología y capacidades productivas. 

	\textcite{Nguyen2020} examinan los determinantes de la complejidad económica en 52 economías, distinguiendo entre países de ingreso alto y medio, e incorporan indicadores de patentes y nueve medidas de desarrollo financiero. Sus estimaciones muestran cointegración de largo plazo y causalidad bidireccional entre patentes, desarrollo financiero y complejidad económica; además, encuentran que las patentes, especialmente las de residentes, se asocian positivamente con la sofisticación productiva, mientras que el efecto del desarrollo financiero depende más de la eficiencia de los mercados que del tamaño del sector financiero.

	\textcite{Sweet2019} analizan el efecto de los derechos de patente sobre el crecimiento de la productividad total de los factores en 70 países durante 1965--2009 mediante regresiones dinámicas de panel. Sus resultados muestran que sistemas de patentes más rigurosos no tienen un efecto significativo sobre la productividad, mientras que la complejidad económica presenta una relación positiva y robusta con ella. Desde la teoría de la capacidad de absorción, los autores interpretan que el crecimiento productivo depende menos de la propiedad formal de innovaciones en la frontera tecnológica que de la capacidad de adaptar, replicar y difundir conocimiento dentro de cadenas productivas internacionales.

	Para países africanos, \textcite{KetuNingaye2024} muestran que la estructura sectorial del empleo también importa para la complejidad económica. Sus estimaciones para 27 países durante 1996--2017 indican que una mayor participación del empleo agrícola se asocia negativamente con la complejidad, mientras que el empleo en industria y servicios contribuye a sostener actividades de mayor sofisticación.
	
	El papel de la innovación aparece con claridad en la evidencia presentada por \textcite{Alvarado2023}. Estos autores examinan los vínculos entre rentas de recursos naturales, complejidad económica e innovación, incorporando además capital humano, ingreso y libertades civiles. Sus resultados muestran que la dependencia de recursos naturales se asocia negativamente con la complejidad económica, pero que la innovación tecnológica ayuda a contrarrestar ese efecto al favorecer la diversificación y la acumulación de capacidades productivas.
	
	La relevancia del nivel de ingreso también se observa en la literatura reciente. \textcite{Li2024} analizan economías recientemente industrializadas con modelos de panel lineales y no lineales, e incorporan ingreso, capital humano, inversión extranjera y densidad poblacional. Sus resultados muestran que la complejidad económica reduce las rentas de recursos naturales en el largo plazo, aunque la relación adopta una forma no lineal de U invertida.
	
	En conjunto, esta literatura sugiere que el capital humano, la inversión extranjera, la innovación y el nivel de ingreso operan como condiciones habilitantes para la diversificación productiva y para la acumulación de capacidades asociadas con mayores niveles de complejidad económica \parencite{YaseminYalta2021,Khan2020,Alvarado2023,Li2024}.

	\section{Condiciones macroeconómicas, enfermedad holandesa y concentración exportadora}
	\label{sec:condiciones-macroeconomicas-dutch-disease}
	
	La evidencia empírica sobre condiciones macroeconómicas y transformación estructural externa en economías intensivas en recursos suele vincularse con los mecanismos de enfermedad holandesa. En este enfoque, los auges de precios o de rentas extractivas pueden apreciar el tipo de cambio real, modificar los incentivos relativos entre actividades transables y no transables, y reducir la competitividad de exportaciones no extractivas. Para el modelo de esta tesis, este canal es relevante porque puede traducirse en menor complejidad económica y en mayor concentración de la canasta exportadora, aun cuando el valor total exportado aumente por el dinamismo de pocos productos primarios.

	En esta línea, \textcite{Bahar2018} muestran que los recursos naturales pueden operar como una forma adicional de maldición de recursos al aumentar la concentración exportadora mediante mecanismos compatibles con la enfermedad holandesa. De manera complementaria, \textcite{Harding2016} documentan que los ingresos extractivos tienden a desplazar exportaciones no asociadas a recursos, lo que reduce el espacio para una canasta exportadora más amplia. Estos resultados son especialmente pertinentes para la variable $DIVX=1-HHI$, ya que sugieren que la dependencia extractiva no solo puede afectar la sofisticación de la estructura exportadora, sino también su grado de diversificación.

	Desde una perspectiva centrada en la complejidad económica, \textcite{Hoang2023} analizan economías avanzadas y emergentes durante 1997--2019 mediante estimaciones de regresión cuantílica en panel. Sus resultados muestran que la incertidumbre de política económica tiende a limitar la evolución de la complejidad, mientras que las rentas de recursos naturales ejercen un efecto negativo consistente con la enfermedad holandesa. Además, el impacto de estos factores varía según el nivel inicial de complejidad y entre economías avanzadas y emergentes, lo que refuerza la importancia de considerar heterogeneidad estructural en el análisis empírico.
	
	\section{Desarrollo financiero y financiamiento de la diversificación productiva}
	\label{sec:desarrollo-financiero-diversificacion}

	\textcite{Mayer2023} analizan el caso de Iraq durante 2000--2022 para evaluar si el desarrollo financiero contribuye a elevar la complejidad económica en una economía fuertemente dependiente del petróleo. Mediante estimaciones DOLS, FMOLS y CCR, los autores encuentran que el desarrollo financiero, el crecimiento económico, la inversión extranjera directa y la urbanización se asocian positivamente con el índice de complejidad económica, mientras que las rentas de recursos naturales presentan un efecto predominantemente negativo. 

	\textcite{Arpaci-Ayhan2023} estudia el efecto de la ayuda externa sobre la complejidad económica en 86 países receptores durante 2003--2019. Sus resultados muestran que la ayuda externa se asocia positivamente con las capacidades productivas de los países receptores, aunque los efectos varían cuando se distingue por sector, modalidad de ayuda y nivel de ingreso. La evidencia también señala que la apertura comercial y la inversión extranjera directa inciden en la complejidad, y que la ayuda puede contribuir a la transformación estructural mediante acumulación de capital humano, infraestructura, instituciones y transferencia de conocimiento.

	\textcite{Ndoya2023} estudian el efecto de los ingresos tributarios sobre la complejidad económica en 29 países africanos durante 1995--2018 mediante estimaciones \textit{System GMM}. Sus resultados muestran que una mayor recaudación tributaria se asocia positivamente con la complejidad, al ampliar los recursos públicos disponibles para financiar bienes y actividades de mayor sofisticación. El análisis de mediación indica, además, que este efecto opera parcialmente a través del desarrollo financiero y del gasto público.
	
	\textcite{Valentine2024} examinan si el desarrollo financiero modifica el vínculo entre recursos naturales y complejidad económica. Mediante estimaciones \textit{System GMM} para 100 países en desarrollo durante 1995--2020, encuentran que la dependencia de recursos reduce la complejidad, pero que el desarrollo financiero atenúa ese efecto adverso. Sus pruebas de robustez muestran que este papel mitigador se mantiene con promedios trianuales y especificaciones no lineales.
	
	En África subsahariana, la dependencia de recursos naturales inhibe la transformación estructural, mientras que el desarrollo financiero puede mitigar ese efecto \parencite{Nwani2023}.
	
	En conjunto, esta evidencia sugiere que el desarrollo financiero puede operar como mecanismo de amortiguación frente a los efectos adversos de la dependencia extractiva y, al mismo tiempo, como condición habilitante para la acumulación de capacidades asociadas con estructuras productivas y exportadoras más diversificadas \parencite{Valentine2024,Nieminen2020,Nwani2023,Lashitew2021}.
	
	\textcite{ValverdeCarbonell2023} desarrollan un enfoque empírico basado en complejidad económica que, mediante un algoritmo no supervisado aplicado a matrices de especialización exportadora, estima simultáneamente dos vectores interdependientes: la competitividad minera de los países, aproximada mediante el índice de aptitud minera (\textit{Mining Fitness Index}), y la criticidad de los minerales, aproximada mediante el índice de criticidad de minerales (\textit{Criticality Minerals Index}), a partir de la estructura de diversidad y ubiquidad de las exportaciones.

	\section{Síntesis de la literatura y conexión con el modelo}
	\label{sec:sintesis-literatura-modelo}

	La evidencia revisada muestra que la transformación estructural externa en economías intensivas en recursos debe analizarse en dos dimensiones complementarias. Por un lado, la complejidad económica permite observar la sofisticación de la canasta exportadora y las capacidades productivas reveladas. Por otro, la diversificación exportadora, aproximada mediante $DIVX=1-HHI$, permite evaluar si la inserción externa depende de pocos productos o si se distribuye en una canasta más amplia, una preocupación recurrente en la literatura sobre recursos naturales, concentración exportadora y diversificación \parencite{Harding2016,Bahar2018,Joya2015,Lashitew2021}. La literatura sobre maldición de los recursos, enfermedad holandesa y dependencia extractiva sugiere que ambas dimensiones pueden verse afectadas por la especialización en recursos, aunque la magnitud y el signo del efecto dependen del contexto institucional, macroeconómico y productivo.

	Los estudios revisados permiten formular expectativas sobre canales
	diferenciados, pero sus resultados no se trasladan automáticamente a la
	muestra de la tesis. En el diseño empírico, $M_1$ y $M_2$ son especificaciones
	temáticas paralelas y $M_3$ integra los canales en una especificación completa
	con efectos fijos de país y año. $ECI$ constituye el resultado principal y
	$DIVX$ el contraste complementario; $RENTS$ es una variable explicativa y
	$DRES$ se utiliza exclusivamente para seleccionar la muestra. De forma
	complementaria, \textcite{Destek2023} muestran que una mayor complejidad
	económica puede reducir la dependencia de recursos naturales, dirección
	inversa que forma parte del debate revisado pero no del estimando principal de
	este trabajo.

	\thesisendchap{literatura-empirica}

\end{document}
